\documentclass[../../../include/open-logic-section]{subfiles}

\begin{document}

\olfileid[id]{sth}{cardinals}{cp}
\olsection{Prinsip Cantor}

Ingat kembali \olref[ordinals][vn]{sec}. Kita sedang membahas himpunan
terurut baik dan mengusulkan bahwa akan berguna jika terdapat objek yang
mewakili pengurutan baik. Dengan gagasan ini, kita memperkenalkan ordinal,
lalu menunjukkan dalam
\olref[ordinals][ordtype]{ordtypesworklikeyouwant} bahwa ordinal berperilaku
sebagaimana kita inginkan, yakni:
\[
	\ordtype{A, <} = \ordtype{B, \lessdot} 
	\text{ jika dan hanya jika } \tuple{A, <} \isomorphic \tuple{B, \lessdot}.
\]
Ingat lebih jauh lagi, yakni \olref[sfr][siz][equ]{sec}. Di sana, dengan
bekerja secara naif, kita memperkenalkan gagasan ``ukuran'' suatu himpunan.
Secara khusus, kita mengatakan bahwa dua himpunan berkardinalitas sama,
$\cardeq{A}{B}$, tepat jika terdapat !!a{bijection} $f \colon A \to B$.
Gagasan ini secara intrinsik {lebih sederhana} daripada pengurutan baik:
kita hanya memperhatikan !!{bijection}s, bukan (seperti pada isomorfisme
urutan) apakah !!{bijection}s tersebut ``mempertahankan suatu struktur''.

Semua ini memunculkan suatu gagasan yang jelas. Sebagaimana kita
memperkenalkan objek-objek tertentu, yakni \emph{ordinal}, untuk mengukur
pengurutan baik, kita dapat memperkenalkan objek-objek tertentu, yakni
\emph{kardinal}, untuk mengukur ukuran. Itulah tujuan bab ini.

Sebelum mengatakan apa yang akan menjadi kardinal-kardinal tersebut, kita
perlu menetapkan suatu prinsip yang harus dipenuhinya. Dengan menulis
$\card{X}$ untuk kardinalitas himpunan $X$, kita ingin agar berlaku:
\[
	\card{A} = \card{B} \text{ jika dan hanya jika } \cardeq{A}{B}.
\]
Kita akan menyebutnya \emph{Prinsip Cantor}, sebab Cantor mungkin merupakan
orang pertama yang memahaminya dengan sangat jelas. (Kita akan membahas lebih
lanjut hubungannya dengan \emph{Prinsip Hume} dalam \olref[hp]{sec}.) Jadi,
tujuan kita ialah mendefinisikan $\card{X}$, untuk setiap $X$, sedemikian
sehingga definisi tersebut menghasilkan Prinsip Cantor.

\end{document}
